\section*{Q-Calculus (Quantum Calculus)}
\addcontentsline{toc}{section}{Q-Calculus (Quantum Calculus)}

\begin{tcolorbox}[
  colback=blue!4,
  colframe=blue!45!black,
  title={Coming Soon},
  arc=2pt,
  left=8pt,right=8pt,top=6pt,bottom=6pt
]
\textbf{This section is under development.}

Q-calculus, also called \emph{quantum calculus}, is a version of calculus built from
\emph{multiplicative} changes rather than additive ones. It replaces the usual derivative
and integral with \(q\)-difference and \(q\)-integral operators, creating a bridge between
continuous calculus, finite calculus, combinatorics, and special functions.

\medskip

Planned topics include:
\begin{itemize}
    \item the \(q\)-difference operator,
    \item \(q\)-numbers and \(q\)-factorials,
    \item the \(q\)-binomial theorem,
    \item the Jackson \(q\)-integral,
    \item basic examples and identities,
    \item the relationship between ordinary calculus, finite calculus, and \(q\)-calculus.
\end{itemize}

\medskip

For now, you may think of q-calculus as a third calculus system in which
\[
\text{derivative} \quad \longleftrightarrow \quad \text{difference} \quad \longleftrightarrow \quad q\text{-difference}.
\]

Like finite calculus, it reveals that many familiar ideas from analysis survive in a new algebraic form.
\end{tcolorbox}

\begin{remark*}
This topic is not needed for the current development, but it will eventually provide a broader
structural view of how different calculi encode different notions of change.
\end{remark*}